\documentclass[11pt]{article}
\usepackage[margin=1in]{geometry}\usepackage{amsmath,booktabs,hyperref}
\title{Where Does Order-Flow-Auction Price Improvement Come From?\\\large Solver Frontiers, Competition, and Batch Selection}
\author{Working-paper implementation}\date{31 August 2026}
\begin{document}\maketitle
\begin{abstract} We combine every submitted CoW solver proposal with contemporaneous archival AMM routes to decompose realised price improvement into a solver frontier and a batch-selection loss.  The central result is an accounting exercise conditional on submitted bids, not a causal claim about entry.  It identifies where solvers create value relative to public liquidity and where the allocation passes or withholds that value from users.\end{abstract}
\section{Introduction} Existing order-flow-auction studies observe outcomes relative to a routed benchmark.  This paper observes the entire proposed execution distribution behind each outcome.  The resulting decomposition distinguishes valuable routing capability from the mechanism that selects among proposed solutions.
\section{Data} The sample contains the full solver-competition census and cached block-pinned Uniswap Smart Order Router quotes.  To preserve a common welfare convention, the automated baseline implementation reports sell orders only; exact-output buy orders require an input-cost transformation and are left for the robustness extension.
\section{Decomposition} For sell order $i$, let $A_i$ be its AMM output, $W_i$ the realised output, and $B_i$ the highest submitted output for the same input.  In basis points relative to notional, $PI_i=(W_i-A_i)/N_i$, $Frontier_i=(B_i-A_i)/N_i$, and $Selection_i=(B_i-W_i)/N_i$.  Hence $PI_i=Frontier_i-Selection_i$.  The identity is descriptive and holds bids fixed.
\section{Results} \texttt{analysis.py} writes the tier-specific sample counts and medians to \texttt{results/results.json}.  The pre-existing full-route analysis already shows substantial median net AMM-relative improvement in mid- and long-tail orders (27.52 and 20.40 bps respectively), but essentially none in the short tail (0.55 bps).  The decomposition determines whether that gradient comes from solver capability or from selection.
\section{Implications} The relevant best-execution question is not simply whether an auction beats an AMM.  It is whether solver innovation creates an attainable frontier, whether competition disciplines it, and whether the batch mechanism transfers it to the order.  This framework makes all three margins auditable.
\section{Conclusion} Full bid data turns price improvement from an outcome statistic into an execution-production account.  The next robustness step is a buy-order analogue, gas-net results, and an exposure event study.
\begin{thebibliography}{9}\bibitem{bachu} Bachu, B., Wan, X., and Moallemi, C. (2024). \emph{Quantifying Price Improvement in Order Flow Auctions}. arXiv:2405.00537.\end{thebibliography}
\end{document}
